% Generated draft: requires mathematical and source review.
\documentclass[11pt]{article}
\usepackage[margin=25mm]{geometry}
\usepackage{fontspec}
\setmainfont{Latin Modern Roman}
\usepackage{amsmath,amssymb,mathtools,unicode-math}
\setmathfont{Latin Modern Math}
\usepackage{enumitem,needspace,xurl,hyperref}
\setlength{\parindent}{0pt}
\setlength{\parskip}{5pt}
\setlength{\emergencystretch}{3em}
\newcommand{\N}{\mathbb{N}}
\newcommand{\Z}{\mathbb{Z}}
\newcommand{\Q}{\mathbb{Q}}
\newcommand{\R}{\mathbb{R}}
\newcommand{\C}{\mathbb{C}}
\newcommand{\F}{\mathbb{F}}
\newcommand{\A}{\mathbb{A}}
\newcommand{\PP}{\mathbb{P}}
\newcommand{\E}{\mathbb{E}}
\newcommand{\eps}{\varepsilon}
\newcommand{\dd}{\,\mathrm{d}}
\begin{document}
\section*{UnsolvedMath ID: 30001721}
\textbf{Source problem code:} catalogue entry
\subsection{Definitions and mathematical statement}
The source asks whether every root of an acyclic quiver admits an indecomposable tree representation. It does not specify the ground field, finiteness hypotheses on the quiver, or the precise root convention. The following gives a parameterized formulation in the customary finite-quiver, positive-root setting; these scope choices require confirmation against the source.

Let \(k\) be a field and let \(Q=(Q_0,Q_1,s,t)\) be a finite quiver with no directed cycles, where \(s,t:Q_1\to Q_0\) give the source and target of each arrow. For distinct vertices \(i,j\), let \(m_{ij}\) be the total number of arrows between \(i\) and \(j\), counting both directions and all parallel arrows. Define the symmetric generalized Cartan matrix \(C_Q\) by
\[
(C_Q)_{ii}=2,\qquad (C_Q)_{ij}=-m_{ij}\quad(i\ne j).
\]
Under the positive-root interpretation, a root dimension vector means a positive root of the Kac--Moody root system associated with \(C_Q\), expressed in the simple-root coordinates indexed by \(Q_0\). Write its set of positive roots as \(\Delta_+(Q)\subseteq\mathbb Z_{\ge0}^{Q_0}\setminus\{0\}\). Thus zero coordinates are allowed, but the zero vector is excluded.

A finite-dimensional representation \(M\) of \(Q\) over \(k\) consists of finite-dimensional \(k\)-vector spaces \(M_i\) for \(i\in Q_0\) and linear maps \(M_a:M_{s(a)}\to M_{t(a)}\) for \(a\in Q_1\). Its dimension vector is \(\underline{\dim}M=(\dim_k M_i)_{i\in Q_0}\). It is indecomposable if it is nonzero and is not isomorphic to a direct sum of two nonzero representations.

Choose a basis \(B_i\) of each \(M_i\), and set \(B=\bigsqcup_{i\in Q_0}B_i\). For \(a:i\to j\) and \(b\in B_i\), write
\[
M_a(b)=\sum_{c\in B_j}\lambda_{a,c,b}c.
\]
The coefficient quiver \(\Gamma(M,B)\) has vertex set \(B\) and one arrow \(b\to c\), labelled by \(a\), for every triple \((a,b,c)\) with \(\lambda_{a,c,b}\ne0\). Call \(M\) a tree representation if some such choice of bases makes the underlying undirected multigraph of \(\Gamma(M,B)\) a tree: it is connected and has no cycles, with distinct coefficient arrows retained as distinct edges. The one-vertex graph with no edges counts as a tree.

With \(k\) retained as an unspecified parameter, the question in this interpretation is whether, for every finite acyclic quiver \(Q\) and every \(\mathbf d\in\Delta_+(Q)\), there exist an indecomposable representation \(M\) over \(k\) and vertex bases \(B\) such that
\[
\underline{\dim}M=\mathbf d
\quad\text{and}\quad
\Gamma(M,B)\text{ has underlying undirected multigraph a tree}.
\]
The record does not determine whether \(k\) is fixed, algebraically closed, arbitrary, or otherwise restricted. Consequently, the displayed question is a parameterized interpretation and does not impose a universal quantifier over fields on the original source.
\subsection{Short English statement}
For each root dimension vector of an acyclic quiver, must there be an indecomposable representation with vertex bases whose nonzero arrow-matrix entries form a tree? The ground field, quiver finiteness, and root convention need confirmation. The supplied record labels the problem partially solved; that label is supplied curation, not an independently verified conclusion.
\subsection{Sources}
\begin{itemize}
\item Supplied local registry record, permanent UnsolvedMath ID \(30001721\), title \emph{Tree Modules for Roots of Acyclic Quivers}, exact source problem code \texttt{catalogue entry}; supplied problem year: 2011. No proposer is named in the proposer field. The supplied literature assessment describes the general question as Ringel's question.
\item Supplied source citation: \emph{Representation Theory of Quivers and Finite Dimensional Algebras} (2011), Oberwolfach Reports, pp.~523--608, DOI 10.4171/owr/2011/10. Supplied literature-source label: ``Oberwolfach Report 10/2011.'' \url{https://doi.org/10.4171/owr/2011/10}. The background describes Oberwolfach Reports as documenting mathematical workshops and their open-problem discussions.
\item Supplied literature source: Weist, \emph{Tree modules}. \url{https://arxiv.org/abs/1011.1203}.
\item Supplied literature source: Weist, ``recursive construction.'' \url{https://www.sciencedirect.com/science/article/pii/S0021869315003737}.
\item Supplied literature assessment, checked \texttt{2026-08-21}: status ``partially solved.'' The supplied assessment states that Weist proves the existence of indecomposable tree modules for every root of the generalized Kronecker quiver and for all imaginary Schur roots, while the question for every root of every acyclic quiver is not proved in the cited generality. This is supplied dated curation derived from the Oberwolfach Reports curation, not independently verified here; the time-sensitive status and cited claims require independent verification.
\end{itemize}
\end{document}
